% GENERATED FILE. Edit canonical lessons, then run npm run generate:books.
% canonical-source-hash: fnv1a64:98a6e8111c53e80c
% canonical-lessons: AR-C345-tahadda, AR-C345-wajaha, AR-C345-taghallaba, AR-C345-iqtaraba, AR-C345-ibtaada

\chapter{Facing danger}
\label{ch:ar-facing-danger}
\hlchaptermodality{345}

\begin{chapteropening}
\textbf{By the end of this chapter:} \emph{I can say to challenge, to face, to overcome, to approach and to move away.}

Complete the last lesson of chapter 345: I can say to challenge, to face, to overcome, to approach and to move away.
\end{chapteropening}

\section[\texorpdfstring{\ar{تحدى} (taḥaddā)}{تحدى (taḥaddā)}]{\ar{تحدى} (taḥaddā) — to challenge, to defy}

\label{lesson:AR-C345-tahadda}



\begin{quote}
Before the new one: say the Arabic for to come true, then the Arabic for to take revenge.
\end{quote}

\subsection*{You'll want to know: \ar{تحدى}}
\textbf{\ar{تحدى}} — \emph{taḥaddā} — \textquotedblleft{}to challenge, to defy\textquotedblright{}.

\begin{quote}
\ar{تحدى} \ar{الولد} \ar{صديقه} \ar{في} \ar{السباق}.
\end{quote}

\begin{quote}
\emph{taḥaddā al-walad ṣadīqahu fī as-sibāq.}
\end{quote}

\begin{quote}
The boy challenged his friend to a race.
\end{quote}

Present: \ar{يتحدى}, \emph{yataḥaddā}.

A challenge is \ar{تحد}; plural \ar{تحديات}.

\begin{cousinweb}
\textbf{\ar{تحدى}} (\emph{taḥaddā}) is built on the native Arabic root \emph{ḥ-d-w} (Lane), where \emph{ḥadā} is to drive camels on with song, and Form V \emph{taḥaddā} is \textquotedblleft{}to vie with someone, to dare him\textquotedblright{}.
\end{cousinweb}

\begin{grammarlens}[title={What you've built}]
One more word for telling what happened.
\end{grammarlens}

\subsection*{Your turn}
\begin{itemize}
\raggedright
  \item \emph{Say it:} \emph{taḥaddā}
  \item \emph{Say it:} \emph{taḥaddā}, once more
  \item \emph{Say it:} \emph{taḥaddā al-walad ṣadīqahu fī as-sibāq.} — The boy challenged his friend to a race.
  \item \emph{Recall:} say the Arabic for to come true, then the Arabic for to take revenge, then say \emph{taḥaddā} again
\end{itemize}

\begin{tcolorbox}[breakable,colback=teal!4,colframe=teal!35!black,title={Before you move on}]
What does \ar{تحدى} mean? (\textquotedblleft{}To challenge, to defy\textquotedblright{}.) Say it once more.
\end{tcolorbox}
\section[\texorpdfstring{\ar{واجه} (wājaha)}{واجه (wājaha)}]{\ar{واجه} (wājaha) — to confront, to face}

\label{lesson:AR-C345-wajaha}



\begin{quote}
Before the new one: say the Arabic for to take revenge, then the Arabic for to challenge.
\end{quote}

\subsection*{You'll want to know: \ar{واجه}}
\textbf{\ar{واجه}} — \emph{wājaha} — \textquotedblleft{}to confront, to face\textquotedblright{}.

\begin{quote}
\ar{واجه} \ar{الفريق} \ar{مشكلة} \ar{كبيرة}.
\end{quote}

\begin{quote}
\emph{wājaha al-farīq mushkila kabīra.}
\end{quote}

\begin{quote}
The team faced a big problem.
\end{quote}

Present: \ar{يواجه}, \emph{yuwājih}.

\begin{cousinweb}
\textbf{\ar{واجه}} (\emph{wājaha}) is built on the native Arabic root \emph{w-j-h} (Lane), the root of the taught \emph{wajh}, a face. It is Form III, \textquotedblleft{}to meet face to face\textquotedblright{}.
\end{cousinweb}

\begin{grammarlens}[title={What you've built}]
One more word for telling what happened.
\end{grammarlens}

\subsection*{Your turn}
\begin{itemize}
\raggedright
  \item \emph{Say it:} \emph{wājaha}
  \item \emph{Say it:} \emph{wājaha}, once more
  \item \emph{Say it:} \emph{wājaha al-farīq mushkila kabīra.} — The team faced a big problem.
  \item \emph{Recall:} say the Arabic for to take revenge, then the Arabic for to challenge, then say \emph{wājaha} again
\end{itemize}

\begin{tcolorbox}[breakable,colback=teal!4,colframe=teal!35!black,title={Before you move on}]
What does \ar{واجه} mean? (\textquotedblleft{}To confront, to face\textquotedblright{}.) Say it once more.
\end{tcolorbox}
\section[\texorpdfstring{\ar{تغلب} (taghallaba)}{تغلب (taghallaba)}]{\ar{تغلب} (taghallaba) — to overcome}

\label{lesson:AR-C345-taghallaba}



\begin{quote}
Before the new one: say the Arabic for to challenge, then the Arabic for to confront.
\end{quote}

\subsection*{You'll want to know: \ar{تغلب}}
\textbf{\ar{تغلب}} — \emph{taghallaba} — \textquotedblleft{}to overcome\textquotedblright{}.

\begin{quote}
\ar{تغلبت} \ar{البنت} \ar{على} \ar{خوفها} \ar{من} \ar{الماء}.
\end{quote}

\begin{quote}
\emph{taghallabat al-bint ʿalā khawfihā min al-māʾ.}
\end{quote}

\begin{quote}
The girl overcame her fear of water.
\end{quote}

Present: \ar{يتغلب}, \emph{yataghallab}.

It takes \ar{على}: \ar{تغلب} \ar{على}.

\begin{cousinweb}
\textbf{\ar{تغلب}} (\emph{taghallaba}) is built on the native Arabic root \emph{gh-l-b} (Lane), \textquotedblleft{}to overpower\textquotedblright{}, the root of the taught \emph{ghāliban}, mostly. It is Form V, \textquotedblleft{}to get the upper hand\textquotedblright{}.
\end{cousinweb}

\begin{grammarlens}[title={What you've built}]
One more word for telling what happened.
\end{grammarlens}

\subsection*{Your turn}
\begin{itemize}
\raggedright
  \item \emph{Say it:} \emph{taghallaba}
  \item \emph{Say it:} \emph{taghallaba}, once more
  \item \emph{Say it:} \emph{taghallabat al-bint ʿalā khawfihā min al-māʾ.} — The girl overcame her fear of water.
  \item \emph{Recall:} say the Arabic for to challenge, then the Arabic for to confront, then say \emph{taghallaba} again
\end{itemize}

\begin{tcolorbox}[breakable,colback=teal!4,colframe=teal!35!black,title={Before you move on}]
What does \ar{تغلب} mean? (\textquotedblleft{}To overcome\textquotedblright{}.) Say it once more.
\end{tcolorbox}
\section[\texorpdfstring{\ar{اقترب} (iqtaraba)}{اقترب (iqtaraba)}]{\ar{اقترب} (iqtaraba) — to approach, to come near}

\label{lesson:AR-C345-iqtaraba}



\begin{quote}
Before the new one: say the Arabic for to confront, then the Arabic for to overcome.
\end{quote}

\subsection*{You'll want to know: \ar{اقترب}}
\textbf{\ar{اقترب}} — \emph{iqtaraba} — \textquotedblleft{}to approach, to come near\textquotedblright{}.

\begin{quote}
\ar{اقترب} \ar{الكلب} \ar{من} \ar{الطفل} \ar{ببطء}.
\end{quote}

\begin{quote}
\emph{iqtaraba al-kalb min aṭ-ṭifl bi-buṭʾ.}
\end{quote}

\begin{quote}
The dog came slowly closer to the child.
\end{quote}

Present: \ar{يقترب}, \emph{yaqtarib}.

It takes \ar{من}: \ar{اقترب} \ar{من}.

\begin{cousinweb}
\textbf{\ar{اقترب}} (\emph{iqtaraba}) is built on the native Arabic root \emph{q-r-b} (Lane), \textquotedblleft{}to be near\textquotedblright{}, the root of the taught \emph{qarīb}. It is Form VIII, \textquotedblleft{}to draw near\textquotedblright{}.
\end{cousinweb}

\begin{grammarlens}[title={What you've built}]
One more word for telling what happened.
\end{grammarlens}

\subsection*{Your turn}
\begin{itemize}
\raggedright
  \item \emph{Say it:} \emph{iqtaraba}
  \item \emph{Say it:} \emph{iqtaraba}, once more
  \item \emph{Say it:} \emph{iqtaraba al-kalb min aṭ-ṭifl bi-buṭʾ.} — The dog came slowly closer to the child.
  \item \emph{Recall:} say the Arabic for to confront, then the Arabic for to overcome, then say \emph{iqtaraba} again
\end{itemize}

\begin{tcolorbox}[breakable,colback=teal!4,colframe=teal!35!black,title={Before you move on}]
What does \ar{اقترب} mean? (\textquotedblleft{}To approach, to come near\textquotedblright{}.) Say it once more.
\end{tcolorbox}
\section[\texorpdfstring{\ar{ابتعد} (ibtaʿada)}{ابتعد (ibtaʿada)}]{\ar{ابتعد} (ibtaʿada) — to move away, to keep away}

\label{lesson:AR-C345-ibtaada}



\begin{quote}
Before the new one: say the Arabic for to overcome, then the Arabic for to approach.
\end{quote}

\subsection*{You'll want to know: \ar{ابتعد}}
\textbf{\ar{ابتعد}} — \emph{ibtaʿada} — \textquotedblleft{}to move away, to keep away\textquotedblright{}.

\begin{quote}
\ar{ابتعد} \ar{الناس} \ar{عن} \ar{البيت} \ar{المحترق}.
\end{quote}

\begin{quote}
\emph{ibtaʿada an-nās ʿan al-bayt al-muḥtariq.}
\end{quote}

\begin{quote}
People moved away from the burning house.
\end{quote}

Present: \ar{يبتعد}, \emph{yabtaʿid}.

It takes \ar{عن}: \ar{ابتعد} \ar{عن}.

\begin{cousinweb}
\textbf{\ar{ابتعد}} (\emph{ibtaʿada}) is built on the native Arabic root \emph{b-ʿ-d} (Lane), \textquotedblleft{}to be far\textquotedblright{}, the root of the taught \emph{baʿīd}. It is Form VIII, \textquotedblleft{}to go far off\textquotedblright{}.
\end{cousinweb}

\begin{grammarlens}[title={What you've built}]
One more word for telling what happened.
\end{grammarlens}

\subsection*{Your turn}
\begin{itemize}
\raggedright
  \item \emph{Say it:} \emph{ibtaʿada}
  \item \emph{Say it:} \emph{ibtaʿada}, once more
  \item \emph{Say it:} \emph{ibtaʿada an-nās ʿan al-bayt al-muḥtariq.} — People moved away from the burning house.
  \item \emph{Recall:} say the Arabic for to overcome, then the Arabic for to approach, then say \emph{ibtaʿada} again
\end{itemize}

\begin{tcolorbox}[breakable,colback=teal!4,colframe=teal!35!black,title={Before you move on}]
What does \ar{ابتعد} mean? (\textquotedblleft{}To move away, to keep away\textquotedblright{}.) Say it once more.
\end{tcolorbox}
